% !TeX root = ../main.tex
% ============================================================
%  sections/04a_notation.tex
% ============================================================

%-----------------------------------------------------------------
\subsection{Assumptions}
\label{sec:model:assumptions}

The model rests on the following assumptions, grouped by the aspect of
the problem they govern.

\medskip
\noindent\textit{Data and cost structure.}
\begin{itemize}[leftmargin=2.8em, label={}, itemsep=0.55em]
\item \refstepcounter{assumption}\label{as:det}%
  \textbf{\theassumption.}\enspace
  Demand $d_{rpt}$ and return-rate bounds
  $[\underline{\eta}_{rpt},\overline{\eta}_{rpt}]$ are known with
  certainty over the entire planning horizon $\setT$.
\item \refstepcounter{assumption}\label{as:lin}%
  \textbf{\theassumption.}\enspace
  All unit economic costs (production, remanufacturing, transport,
  inventory holding, and material acquisition) are constant across
  periods and scale linearly with quantities.
\item \refstepcounter{assumption}\label{as:emis}%
  \textbf{\theassumption.}\enspace
  All emission factors are constant and technology- or mode-specific;
  emissions scale linearly with processed output, transported mass,
  and distance.
\item \refstepcounter{assumption}\label{as:lead}%
  \textbf{\theassumption.}\enspace
  Transportation lead times are absorbed within a single planning period.
\end{itemize}

\medskip
\noindent\textit{Operational constraints.}
\begin{itemize}[leftmargin=2.8em, label={}, itemsep=0.55em]
\item \refstepcounter{assumption}\label{as:dem}%
  \textbf{\theassumption.}\enspace
  Every market demand is met exactly in each period; no inventory is
  held at markets.
\item \refstepcounter{assumption}\label{as:hub}%
  \textbf{\theassumption.}\enspace
  Hub terminals operate as pure cross-docking facilities: all items
  arriving at a hub in period $t$ depart in the same period, with no
  on-site storage. Following \citet{mota2018}, hubs are treated as
  outsourced terminal-access points; selecting a hub incurs fixed access
  costs and flow-dependent handling costs, but no facility-construction
  or opening emissions.
\item \refstepcounter{assumption}\label{as:feas}%
  \textbf{\theassumption.}\enspace
  For every admissible index $\taxhat\in\{\taxlb,\ldots,\taxub\}$,
  the follower's feasible region $\setX$ is non-empty. This can be
  verified a priori by checking that the network has sufficient
  installed capacity to cover maximum aggregate demand.
\end{itemize}

\medskip
\noindent\textit{Bilevel game structure.}
\begin{itemize}[leftmargin=2.8em, label={}, itemsep=0.55em]
\item \refstepcounter{assumption}\label{as:stackelberg}%
  \textbf{\theassumption.}\enspace
  The regulator announces the carbon-tax index $\taxhat$ before the
  supply chain operator makes any facility, technology, or operational
  decision; the operator observes $\taxhat$ with certainty and responds
  rationally (\emph{open-loop Stackelberg game}).
\item \refstepcounter{assumption}\label{as:fixednet}%
  \textbf{\theassumption.}\enspace
  Facility-opening, supplier and hub selection, and
  technology-installation decisions are made at the start of the
  planning horizon and are irrevocable; no opening, closure, selection
  change, or capacity expansion occurs mid-horizon.
\item \refstepcounter{assumption}\label{as:singleop}%
  \textbf{\theassumption.}\enspace
  The supply chain operator is a single entity that centrally
  coordinates all network design and operational decisions across all
  echelons.
\item \refstepcounter{assumption}\label{as:opt}%
  \textbf{\theassumption.}\enspace
  When the follower's optimal response is non-unique, the leader
  selects the response with the lowest total emissions
  (\emph{optimistic bilevel convention}).
\end{itemize}

%-----------------------------------------------------------------
\subsection{Notation}
\label{sec:model:notation}

The following notation is used throughout the formulation. All
sets, parameters, and decision variables are summarised in
\Cref{tab:sets,tab:econ_params,tab:capacity_params,tab:demand_params,%
tab:emission_params,tab:variables}.

% !TeX root = ../main.tex
% ============================================================
%  tables/tab_notation.tex
%  Sets, parameters, and follower decision variables.
% ============================================================

%--------------------------------------------------------------
\begin{table}[htbp]
\centering
\caption{Sets and indices.}
\label{tab:sets}
\footnotesize
\renewcommand{\arraystretch}{1.15}
\begin{tabular}{@{}lll@{}}
\toprule
Set & Index & Description \\
\midrule
\multicolumn{3}{@{}l}{Network nodes} \\
$\setS$ & $s$ & Suppliers \\
$\setF$ & $f$ & Factories \\
$\setW$ & $w$ & Warehouses \\
$\setR$ & $r$ & Markets (retailers) \\
$\setHa,\setHs$ & $h$ & Airports; seaports \\
$\setH$ & $h$ & Hub terminals, $\setH\defeq\setHa\cup\setHs$ \\
$\setN$ & $i,j$ & All nodes, $\setN\defeq\setS\cup\setF\cup\setW\cup\setR\cup\setH$ \\
$\setNo$ & $i$ & Openable nodes, $\setNo\defeq\setS\cup\setF\cup\setW\cup\setH$ \\
$\setNv$ & $i$ & Inventory nodes, $\setNv\defeq\setW$ \\
\addlinespace
\multicolumn{3}{@{}l}{Items} \\
$\setU,\setP,\setK$ & $u,p,k$ & Raw materials; finished products; recoverable items \\
$\setQ$ & $q$ & All transportable items, $\setQ\defeq\setU\cup\setP\cup\setK$ \\
\addlinespace
\multicolumn{3}{@{}l}{Technologies} \\
$\setG,\setL$ & $g,\ell$ & Production and remanufacturing technologies \\
$\Omega^{\mathrm{p}}$ & $(p,g)$ & Product-technology pairs for production, $\Omega^{\mathrm{p}}\subseteq\setP\times\setG$ \\
$\Omega^{\mathrm{r}}$ & $(p,\ell)$ & Product-technology pairs for remanufacturing, $\Omega^{\mathrm{r}}\subseteq\setP\times\setL$ \\
\addlinespace
\multicolumn{3}{@{}l}{Transport modes} \\
$\setMrd$ & $m$ & Company-owned trucks (road mode) \\
$\setMa,\setMs$ & $m$ & Air; sea transfer modes \\
$\setMout$ & $m$ & Outsourced modes, $\setMout\defeq\setMa\cup\setMs$ \\
$\setM$ & $m$ & All transport modes, $\setM\defeq\setMrd\cup\setMout$ \\
\addlinespace
\multicolumn{3}{@{}l}{Arcs} \\
$\setA$ & $(i,j)$ & Admissible arcs \\
$\setA_q$ & $(i,j)$ & Admissible arcs for item $q$ \\
$\setM_{ij}$ & $m$ & Modes permitted on arc $(i,j)$ \\
\addlinespace
\multicolumn{3}{@{}l}{Time} \\
$\setT$ & $t$ & Planning periods, $\setT\defeq\{1,\ldots,T\}$ \\
\bottomrule
\end{tabular}
\end{table}

%--------------------------------------------------------------
\begin{table}[htbp]
\centering
\caption{Economic parameters.}
\label{tab:econ_params}
\footnotesize
\renewcommand{\arraystretch}{1.15}
\begin{tabular}{@{}llr@{}}
\toprule
Parameter & Description & Units \\
\midrule
$F_i$ & Fixed cost of selecting supplier $i\in\setS$ & USD \\
$c^{\mathrm{area}}_i$ & Installation cost per unit of floor area at facility $i\in\setF\cup\setW$ & USD/m$^2$ \\
$F_h$ & Per-period access (contract) cost of hub $h\in\setH$ & USD/period \\
$C^{\mathrm{p}}_g,\;C^{\mathrm{r}}_{\ell}$ & Installation cost of production/remanufacturing technology & USD \\
$c^{\mathrm{p}}_g,\;c^{\mathrm{r}}_{\ell}$ & Unit operating cost of production/remanufacturing & USD/unit \\
$c^{\mathrm{u}}_{su}$ & Unit acquisition cost of raw material $u$ from supplier $s$ & USD/unit \\
$c^{\mathrm{k}}_k$ & Unit acquisition cost of recovered item $k$ & USD/unit \\
$c^{\mathrm{h}}_p$ & Unit holding cost of finished product $p$ per period (warehouses only) & USD/(unit$\cdot$period) \\
$c^{\mathrm{tr}}_m$ & Variable outsourced transport cost, $m\in\setMout$ & USD/(kg$\cdot$km) \\
$c^{\mathrm{veh}}_m$ & Road transport cost per vehicle-km, $m\in\setMrd$ & USD/(veh$\cdot$km) \\
$c^{\mathrm{hub}}_m$ & Unit handling cost at a hub terminal for mode $m$ & USD/kg \\
$F^{\mathrm{m}}_m$ & Vehicle investment cost, road mode $m\in\setMrd$ & USD \\
\bottomrule
\end{tabular}
\end{table}

%--------------------------------------------------------------
\begin{table}[htbp]
\centering
\caption{Capacity, resource, and fleet parameters.}
\label{tab:capacity_params}
\footnotesize
\renewcommand{\arraystretch}{1.15}
\begin{tabularx}{\textwidth}{@{}lXr@{}}
\toprule
Parameter & Description & Units \\
\midrule
$\underline{S}_{su}$ & Min. per-period supply of raw material $u$ from supplier $s$ & units/period \\
$\overline{S}_{su}$ & Max. per-period supply of raw material $u$ from supplier $s$ & units/period \\
$\underline{Q}^{\mathrm{p}}_{g}$ & Min. aggregate output of production technology $g$ per period & units/period \\
$\overline{Q}^{\mathrm{p}}_{g}$ & Max. aggregate output of production technology $g$ per period & units/period \\
$\underline{Q}^{\mathrm{r}}_{\ell}$ & Min. aggregate output of remanufacturing technology $\ell$ per period & units/period \\
$\overline{Q}^{\mathrm{r}}_{\ell}$ & Max. aggregate output of remanufacturing technology $\ell$ per period & units/period \\
$\underline{I}_{pi}$ & Min. inventory of finished product $p$ at warehouse $i\in\setNv$ & units \\
$\overline{I}_{pi}$ & Max. inventory of finished product $p$ at warehouse $i\in\setNv$ & units \\
$I^{0}_{pi}$ & Initial inventory ($t=0$) of finished product $p$ at warehouse $i\in\setNv$ & units \\
$\underline{A}_{i}$ & Min. installation area of facility $i\in\setF\cup\setW$ & m$^2$ \\
$\overline{A}_{i}$ & Max. installation area of facility $i\in\setF\cup\setW$ & m$^2$ \\
$\overline{B}_m$ & Contracted per-arc-period capacity of mode $m\in\setMout$ & kg \\
$\underline{W}_{m}$ & Min. arc-period weight when mode $m$ is active & kg \\
$\overline{W}_{m}$ & Max. arc-period weight when mode $m$ is active & kg \\
$M^{\mathrm{trip}}_m$ & Big-$M$ bound on trips by mode $m$ per arc-period & -- \\
$r_m$ & Return-distance factor ($r_m=1$ for outsourced modes) & -- \\
$v_m,H_m$ & Average speed, driving-hour limit (road) & km/h, h \\
$\overline{V}$ & Max. total investment in company-owned vehicles & USD \\
$apu_q$ & Storage area per unit of item $q\in\setQ$ held in end-of-period inventory; only $q\in\setP$ at warehouses carries nonzero stock~\eqref{eq:YCT}, since raw materials $\setU$ and recovered items $\setK$ hold no on-site inventory state anywhere & m$^2$/unit \\
$apur_p$ & Handling area per unit of period throughput of finished product $p\in\setP$, entering the flow-based component of~\eqref{eq:YCT} at both factories and warehouses & m$^2$/(unit/period) \\
\bottomrule
\end{tabularx}
\end{table}

%--------------------------------------------------------------
\begin{table}[htbp]
\centering
\caption{Demand, return, and conversion parameters.}
\label{tab:demand_params}
\footnotesize
\renewcommand{\arraystretch}{1.15}
\begin{tabular}{@{}llr@{}}
\toprule
Parameter & Description & Units \\
\midrule
$\alpha_{upg}$ & Raw-material input ratio for $u$ per unit of $p$ via $g$ & unit/unit \\
$\beta_{kp}$ & Recovered-item input ratio for $k$ per remanufactured unit of $p$ & unit/unit \\
$\gamma_{kp}$ & Yield of recovered item $k$ per returned unit of $p$ & unit/unit \\
$\omega_q$ & Unit weight of item $q$ & kg \\
$d_{rpt}$ & Demand for product $p$ at market $r$ in period $t$ & units/period \\
$\underline{\eta}_{rpt}$ & Min. return-rate fraction for $p$ at market $r$ in period $t$ & -- \\
$\overline{\eta}_{rpt}$ & Max. return-rate fraction for $p$ at market $r$ in period $t$ & -- \\
$\theta_p$ & Integer return lag of product $p$ & periods \\
$\delta_{ij}$ & Distance along arc $(i,j)$ & km \\
\bottomrule
\end{tabular}
\end{table}

%--------------------------------------------------------------
\begin{table}[htbp]
\centering
\caption{Emission-intensity and leader-policy parameters.}
\label{tab:emission_params}
\footnotesize
\renewcommand{\arraystretch}{1.15}
\begin{tabular}{@{}llr@{}}
\toprule
Parameter & Description & Units \\
\midrule
$e^{\mathrm{o,f}}_i$ & Opening emission intensity for facility $i\in\setF\cup\setW$ & kg\,CO$_2$e/m$^2$ \\
$e^{\mathrm{p}}_g,\;e^{\mathrm{r}}_{\ell}$ & Emission intensity per kg produced/remanufactured & kg\,CO$_2$e/kg \\
$e^{\mathrm{u}}_{su}$ & Emission intensity of raw material $u$ from supplier $s$ & kg\,CO$_2$e/kg \\
$e^{\mathrm{tr}}_m$ & Transport emission intensity of mode $m$ & kg\,CO$_2$e/(kg$\cdot$km) \\
$\taxindexlb,\taxindexub$ & Lower/upper bounds on the leader's carbon-tax index & -- \\
\bottomrule
\end{tabular}
\end{table}

%--------------------------------------------------------------
\begin{table}[htbp]
\centering
\caption{Follower decision variables.}
\label{tab:variables}
\footnotesize
\renewcommand{\arraystretch}{1.15}
\begin{tabular}{@{}lll@{}}
\toprule
Variable & Domain & Description \\
\midrule
$y_i$ & $\bin$ & 1 if node $i\in\setNo$ is activated (opened, if $i\in\setF\cup\setW$; selected, if $i\in\setS\cup\setH$) \\
$YC_i$ & $\Rplus$ & Installed floor area at facility $i\in\setF\cup\setW$ (m$^2$) \\
$\lambda^{\mathrm{p}}_{fg}$ & $\bin$ & 1 if production technology $g$ is installed at $f$ (shared across products) \\
$\lambda^{\mathrm{r}}_{f\ell}$ & $\bin$ & 1 if remanufacturing technology $\ell$ is installed at $f$ (shared across products) \\
$a_{f\ell t}$ & $\bin$ & 1 if installed technology $\ell$ is operated at $f$ in period $t$ \\
$\sigma_{sut}$ & $\bin$ & 1 if supplier $s$ ships raw material $u$ in period $t$ \\
$\zeta_{ijmt}$ & $\bin$ & 1 if mode $m$ is active on arc $(i,j)$ in period $t$ \\
$\kappa_{mt}$ & $\bin$ & 1 if outsourced capacity for mode $m$ is contracted in period $t$ \\
$Q_{ijmt}$ & $\Zplus$ & Trips on arc $(i,j)$ by road mode $m\in\setMrd$ in period $t$ \\
$K_{im}$ & $\Zplus$ & Company-owned vehicles of road mode $m$ assigned to node $i$ \\
$q^{\mathrm{p}}_{pfgt}$ & $\Rplus$ & Production of $p$ at $f$ via technology $g$ in $t$, $(p,g)\in\Omega^{\mathrm{p}}$ \\
$q^{\mathrm{r}}_{pf\ell t}$ & $\Rplus$ & Remanufacturing of $p$ at $f$ via $\ell$ in $t$, $(p,\ell)\in\Omega^{\mathrm{r}}$ \\
$s^{\mathrm{inv}}_{pit}$ & $\Rplus$ & Inventory of finished product $p$ at warehouse $i\in\setNv$ in period $t$ \\
$x_{qijmt}$ & $\Rplus$ & Flow of item $q$ on arc $(i,j)$ by mode $m$ in period $t$ \\
\bottomrule
\end{tabular}
\end{table}


%- - - - - - - - - - - - - - - - - - - - - - - - - - - - - - - -
\subsubsection{Sets and indices}
\label{sec:model:sets}

\Cref{tab:sets} lists the node, item, technology, transport, arc,
and planning-period sets. The set of all transportable items is
$\setQ=\setU\cup\setP\cup\setK$ and the set of all nodes is
$\setN=\setS\cup\setF\cup\setW\cup\setR\cup\setH$.
Openable nodes are $\setNo=\setS\cup\setF\cup\setW\cup\setH$;
inventory-holding nodes are $\setNv=\setW$ -- factories carry no
inventory of their own (\Cref{sec:model:c1}), so finished-product
stock is held only at warehouses.
Transport partition as $\setM=\setMrd\cup\setMout$, where
$\setMout=\setMa\cup\setMs$ collects the outsourced non-road modes.
Instance files additionally carry a set labelled $I$, kept equal to
$\setN$ for backward compatibility with an earlier node-indexing
convention; it is validated against $\setN$ on load but is not an
independent model object, and the implementation reads only $\setN$.

%- - - - - - - - - - - - - - - - - - - - - - - - - - - - - - - -
\subsubsection{Parameters}
\label{sec:model:parameters}

\Cref{tab:econ_params} lists economic cost parameters;
\Cref{tab:capacity_params} lists capacity, resource, and fleet
parameters; \Cref{tab:demand_params} lists demand, return, and
conversion parameters; \Cref{tab:emission_params} lists emission
intensities and the leader's policy bounds. The set
$\setM_{ij}\subseteq\setM$ denotes the modes permitted on arc $(i,j)$,
and $\delta_{ij}$ its distance. $\setM_{ij}$ is general per-arc
instance data; the data generator and implementation realise it via a
fixed topological rule -- air-air arcs admit only $\setM^{\mathrm{a}}$,
sea-sea arcs only $\setM^{\mathrm{s}}$, and every other arc admits the
road-family modes $\setMrd$ -- which is sufficient for every benchmark
topology used in this paper but is not the most general case the
notation $\setM_{ij}$ permits. The indicator
\begin{equation*}
\xi^{\mathrm{o}}_{ij}=\mathbf{1}[i\in\setH], \qquad
\xi^{\mathrm{d}}_{ij}=\mathbf{1}[j\in\setH]
\end{equation*}
flags whether the origin or destination of arc $(i,j)$ is a hub
terminal, used to assess handling charges. Compatible
product--technology pairs are collected in
$\Omega^{\mathrm{p}}\subseteq\setP\times\setG$ (production) and
$\Omega^{\mathrm{r}}\subseteq\setP\times\setL$ (remanufacturing).

%- - - - - - - - - - - - - - - - - - - - - - - - - - - - - - - -
\subsubsection{Decision variables}
\label{sec:model:variables}

\paragraph{Leader decision variable.}
The regulator's sole decision is the integer carbon-tax index
\begin{equation}
\taxhat\in\mathcal{T}^{\tau}\defeq
\{\taxlb,\taxlb+1,\ldots,\taxub\}\subset\Z,
\label{eq:leader_domain}
\end{equation}
which induces the carbon-tax rate $\tau=\taxhat/100$ (USD/kg\,CO$_2$e).
We work with the index $\taxhat$ rather than the rate $\tau$ directly
because it keeps the leader's domain finite and discrete, enabling 
the exact DBP of \Cref{sec:methodology}.

\paragraph{Follower decision variables.}
Given $\taxhat$, the follower chooses the vector
\begin{equation*}
\bz=\bigl(y,\,YC,\,\lambda^{\mathrm{p}},\,\lambda^{\mathrm{r}},\,a,\,
\sigma,\,\zeta,\,\kappa,\,Q,\,K,\,q^{\mathrm{p}},\,q^{\mathrm{r}},\,
s^{\mathrm{inv}},\,x\bigr).
\end{equation*}
\Cref{tab:variables} lists all component variables and their domains.

To decouple the algebraic model from the physical network, the
admissible flow-tuple set is
\begin{equation}
\Phi \defeq
\bigl\{(q,i,j)\,:\, q\in\setQ,\,(i,j)\in\setA_q\bigr\},
\label{eq:phi}
\end{equation}
which admits the partition
$\Phi=\Phi^{\mathrm{u}}\cup\Phi^{\mathrm{p}}\cup\Phi^{\mathrm{k}}$
according to $q\in\setU$, $q\in\setP$, or $q\in\setK$.

%- - - - - - - - - - - - - - - - - - - - - - - - - - - - - - - -
\subsubsection{Auxiliary expressions}
\label{sec:model:auxiliary}

The following shorthand expressions are used throughout the
formulation to keep the notation compact. For each item $q\in\setQ$,
node $i\in\setN$, and period $t\in\setT$, the aggregate inflow and
outflow at node $i$ are
\begin{equation}
X^{\mathrm{in}}_{qit} \defeq
  \sum_{j:(j,i)\in\setA_q}\,\sum_{m\in\setM_{ji}} x_{qjimt},
\qquad
X^{\mathrm{out}}_{qit} \defeq
  \sum_{j:(i,j)\in\setA_q}\,\sum_{m\in\setM_{ij}} x_{qijmt}.
\label{eq:flows}
\end{equation}
To avoid case distinctions between openable and non-openable nodes
in the flow constraints,
\begin{equation}
\bar{y}_{i} \defeq
\begin{cases}
y_{i}, & i\in\setNo,\\
1,     & i\in\setR.
\end{cases}
\label{eq:ybar}
\end{equation}
The used floor area at facility $i\in\setF\cup\setW$ in period $t$
combines a flow-based handling component, present at both factories
and warehouses, with a stock-based storage component, present only at
warehouses (factories carry no inventory of any kind,
\Cref{sec:model:c1}):
\begin{equation}
YCT_{it} \defeq
\begin{cases}
\displaystyle\sum_{p\in\setP} apur_p\,X^{\mathrm{out}}_{pit},
  & i\in\setF,\\[1.1em]
\displaystyle
\sum_{p\in\setP} apur_p\,X^{\mathrm{in}}_{pit}
  + \sum_{p\in\setP} apu_p\,s^{\mathrm{inv}}_{pit},
  & i\in\setW,
\end{cases}
\qquad t\in\setT.
\label{eq:YCT}
\end{equation}
Factories carry no inventory of any kind (\Cref{sec:model:c1}), so
their floor-area footprint has no stock component; their handling
space is driven entirely by per-period throughput. Warehouses carry
both: incoming-flow handling space and finished-product stock space,
following \citet{mota2018}'s general facility-area expression (their
Eq.~16), adapted to this project's flow-direction convention --- see
\Cref{rem:floor_area_flow_direction}.
